%! Semester 1, Units 1–3 — Final Exam Study Guide (reference sheet; no answer key by design)
% This is a REFERENCE sheet, not a worksheet: every worked move and every trap is already
% filled in, so there is nothing to answer and therefore no _key sibling. The practice a
% student does before the final is ../practice_test/, which HAS a key.
% NAMESTRIP: no \namedateperiod — the name row lives on a packet cover only.
% NOT GRADED: a study guide carries no \gradedblank and does not load grading (LESSON_SHAPE §9).
% It follows the final's blueprint move for move (see ../actual_test/main.tex): the 20
% multiple-choice targets, the five short-answer computations, and the three free-response
% questions. Vocabulary is NOT re-listed here — the unit study guides (unit01/quizzes,
% unit02/tests, unit03/tests) carry it; this sheet is the moves, the sentences, and the traps.
% Worked contexts (a bike-share program; piano practice and recital scores; a community
% garden) are deliberately NOT used on either form of the final.
%
% THE DATA (every number checked)
%  Move 3  12 15 18 20 21 24 26 27 30 33 41: median 24, Q1 18, Q3 30, IQR 12; fences 0 and 48.
%  Move 6  N(50, 6): 44–56 is 68%; above 62 is 2.5%; 59 -> z = 1.5.
%  Move 7  members 84/36 (120), casual 96/144 (240), totals 180/180/360:
%          84/120 = 0.70 of members ride on weekdays; 84/180 = 0.467 of weekday riders are members.
%  Move 9  yhat = 14.20 + 3.150(hours); R-Sq 81.0% -> r = +0.90; 8 hours -> 39.4; scored 42 ->
%          residual 2.6.
%  Move 15 100 reshuffles, step 1: -4:1 -3:2 -2:12 -1:21 0:28 1:21 2:12 3:2 4:1 (sum 100).
%          Observed 3: 2 + 1 = 3 of 100 at or above -> significant.
\documentclass[12pt]{article}
\usepackage{apstats-article}
\usepackage{apstats-boxes}
\usepackage{ltablex}
\keepXColumns

% Optional arg: the \boxguard size (lines kept with the strip).
\newcommand{\sghead}[2][10]{%
  \vspace{0.13in}
  \boxguard[#1]
  \begin{headlinebox}{navy}\color{white}\bfseries\small #2\end{headlinebox}
  \vspace{0.09in}}

\begin{document}

\pageheader{Semester 1 \quad Units 1--3}{Final Exam Study Guide}

\vspace{0.12in}

\boxguard[9]
\begin{remindbox}
\small
The final covers \textbf{Units 1, 2, and 3} --- describing data, relationships between two
variables, and how data are collected. You have \textbf{up to 2 hours}. \textbf{Section I:} 20
multiple-choice questions (40 points). \textbf{Section II:} 5 short answers --- a number in a
box (15 points). \textbf{Section III:} 3 free-response questions, one per unit, five parts each
(45 points). Every written answer must be \textbf{in context} --- name the people or things
studied, the variable, and its units.
\end{remindbox}

\sghead{Unit 1 --- Exploring one-variable data}

\small
\textbf{Move 1 --- Categorical or quantitative? (1.1)} Ask whether the values are \emph{measured
or counted amounts}. Digits do not decide it: a station number, a ZIP code, and a jersey number
are labels, so they are \textbf{categorical}. Trip length in minutes is \textbf{quantitative}.

\vspace{3pt}
\textbf{Move 2 --- Compare groups with proportions, not counts (1.2).} \emph{Worked:} 45 of 150
weekday riders and 40 of 100 weekend riders wear a helmet. Weekday has more helmets (45), but
helmets are \textbf{more common on weekends}: 40\% compared with 30\%.

\vspace{3pt}
\textbf{Move 3 --- Median, quartiles, IQR, and the outlier rule (1.3, 1.5, 1.6).} \emph{Worked,}
11 trip lengths in minutes: 12 \ 15 \ 18 \ 20 \ 21 \ 24 \ 26 \ 27 \ 30 \ 33 \ 41. The median is
the 6th value, \textbf{24}. $Q_1$ is the median of the lower five, \textbf{18}; $Q_3$ is the
median of the upper five, \textbf{30}. $\text{IQR} = 30 - 18 = 12$. Fences:
$18 - 1.5(12) = 0$ and $30 + 1.5(12) = 48$. No trip is below 0 or above 48, so there are
\textbf{no outliers}. From a stem plot, do exactly the same: count in to the middle.

\vspace{3pt}
\textbf{Move 4 --- Shape tells you where the mean is (1.4, 1.5).} The mean is pulled toward the
tail. \textbf{Skewed right:} mean $>$ median. \textbf{Skewed left:} mean $<$ median.
\textbf{Symmetric:} about equal. The median and the IQR are \textbf{resistant}; the mean, the
standard deviation, and the range are not. Change one extreme value and the mean moves, while
the median stays put as long as the value stays on the same side of the middle.

\vspace{3pt}
\textbf{Move 5 --- Read a boxplot (1.5, 1.6).} Each of the four sections holds about
\textbf{a quarter of the data}, however long it is drawn: a long whisker means the values are
spread out, not that there are more of them. A value at the \textbf{70th percentile} has 70\% of
the data at or below it.

\vspace{3pt}
\boxguard[7]
\textbf{Move 6 --- Parameter or statistic; $z$-scores; the empirical rule (1.7).} A
\textbf{parameter} describes the whole population; a \textbf{statistic} describes a sample.
\emph{Worked:} trip lengths are approximately normal with mean 50 and standard deviation 6
minutes. About \textbf{68\%} of trips last 44 to 56 minutes (within 1 SD), \textbf{95\%} last 38
to 62 (within 2), and 99.7\% last 32 to 68 (within 3). Longer than 62 minutes: half of the other
5\%, about \textbf{2.5\%}. A 59-minute trip has $z = \dfrac{59 - 50}{6} = 1.5$: it is 1.5
standard deviations above the mean. To compare values from two different distributions, compare
their $z$-scores.

\normalsize
\sghead{Unit 2 --- Exploring two-variable data}

\small
\boxguard[9]
\textbf{Move 7 --- ``Of the \ldots'' names the denominator (2.1).}

\vspace{2pt}
\noindent\begin{minipage}[c]{0.42\linewidth}\centering\footnotesize
\renewcommand{\arraystretch}{1.15}
\begin{tabular}{@{} l c c c @{}}
\rowcolor{sky}
 & \textbf{Weekday} & \textbf{Weekend} & \textbf{Total} \\
\midrule
Member & 84 & 36  & 120 \\
Casual & 96 & 144 & 240 \\
\midrule
Total  & 180 & 180 & 360 \\
\end{tabular}
\end{minipage}\hfill
\begin{minipage}[c]{0.55\linewidth}\raggedright
\emph{Worked:} of the \textbf{members}, the proportion who rode on a weekday is $84/120 = 0.70$.
Of the \textbf{weekday rides}, the proportion taken by members is $84/180 \approx 0.467$. Same
84, different question. Rider type and day are \textbf{associated}: 70\% of members ride on
weekdays, but only 40\% of casual riders do.
\end{minipage}

\vspace{5pt}
\textbf{Move 8 --- Describe a scatterplot; know what $r$ is (2.2).} Give \textbf{direction, form,
strength}, and any \textbf{unusual points}, in context. The correlation $r$ has no units, does
not change when units change, measures only \emph{linear} association, and never shows cause.

\vspace{3pt}
\boxguard[11]
\textbf{Move 9 --- Read computer output and use the line (2.3, 2.4).}

\vspace{2pt}
\noindent\begin{minipage}[c]{0.40\linewidth}\centering\scriptsize\ttfamily
\begin{tabular}{@{} l r r @{}}
Predictor & Coef & SE Coef \\ \hline
Constant & 14.20 & 2.861 \\
Hours & 3.150 & 0.3596 \\
\end{tabular}\par
\vspace{2pt}
S = 4.11 \quad R-Sq = 81.0\%
\end{minipage}\hfill
\begin{minipage}[c]{0.57\linewidth}\raggedright
\emph{Worked,} recital score from weekly practice hours for 20 piano students:
$\widehat{\text{score}} = 14.20 + 3.150(\text{hours})$. The slope is on the row named for the
explanatory variable. Since the slope is positive, $r = +\sqrt{0.810} = 0.90$. A student who
practices 8 hours is predicted to score $14.20 + 3.150(8) = 39.4$. She scored 42, so her
\textbf{residual} is actual $-$ predicted $= 42 - 39.4 = 2.6$ points.
\end{minipage}

\vspace{5pt}
\textbf{Move 10 --- Is the line appropriate, and where does it stop? (2.3, 2.5)} A residual plot
with \textbf{random scatter} around 0 says a line fits; a \textbf{curve} says it does not, even
when $r$ is high. Predicting far outside the $x$-values in the data is \textbf{extrapolation},
and the answer may be nonsense. A point far from the others in the $x$-direction can be
\textbf{influential}: remove it and the slope and $r$ change a lot.

\normalsize
\sghead{Unit 3 --- Collecting data}

\small
\textbf{Move 11 --- Observational study or experiment? (3.1)} One question: \emph{did the
researchers assign a treatment?} Careful measuring, an app, or a random sample does not make a
study an experiment.

\vspace{3pt}
\textbf{Move 12 --- Choose a sample, and name the method (3.2).} \emph{SRS:} label, randomize
with a random number generator ignoring repeats, select. \textbf{Stratified:} an SRS from
\emph{every} group. \textbf{Cluster:} \emph{all} of a few randomly chosen groups.

\vspace{3pt}
\textbf{Move 13 --- Name the bias, who, and which way (3.3).} \textbf{Voluntary response}
(people choose to answer), \textbf{undercoverage} (some cannot be chosen), \textbf{nonresponse}
(chosen people cannot be reached or refuse), \textbf{response bias} (the question pushes the
answer). Say who is over- or under-represented and whether the estimate is too high or too low.
\textbf{A larger sample does not fix bias.}

\vspace{3pt}
\boxguard[6]
\textbf{Move 14 --- Lay out an experiment (3.4).} \emph{Worked:} a community garden tests 2
kinds of mulch at 3 watering levels on 60 tomato plants. \textbf{Units:} the 60 plants.
\textbf{Treatments:} every combination, $2 \times 3 = 6$. \textbf{Response:} kilograms of
tomatoes per plant. \textbf{Random assignment} balances other variables, such as sunlight, so
they are not \textbf{confounded} with the treatments. If sunlight matters, \textbf{block}: sort
the plants into sunny and shady blocks, then randomly assign \emph{within} each block.

\vspace{3pt}
\boxguard[14]
\textbf{Move 15 --- Read a reshuffle dotplot (3.5).} Each dot is one difference that chance alone
produced.

\vspace{2pt}
\noindent\begin{minipage}[c]{0.44\linewidth}\raggedright
\emph{Worked:} the mulched plants' mean yield was \textbf{3 kg higher}. Count the dots at 3 or
more: $2 + 1 =$ \textbf{3 of 100}. Chance alone rarely gives a difference that large, so it
is \textbf{statistically significant}. Count on the side the observed difference is on, and
include the column it sits on. Had 30 of the 100 dots reached it, chance could easily explain
it: \textbf{not} significant.
\end{minipage}\hfill
\begin{minipage}[c]{0.53\linewidth}\centering
\begin{tikzpicture}[x=0.85cm, y=0.105cm]
  \draw[thick] (-4.5,0) -- (4.5,0);
  \foreach \x in {-4,...,4} {
    \draw (\x,0) -- (\x,-1.6); \node[below, font=\tiny] at (\x,-1.4) {$\x$};}
  \foreach \x/\n in {-4/1,-3/2,-2/12,-1/21,0/28,1/21,2/12,3/2,4/1}
    \foreach \k in {1,...,\n} \fill[navy] (\x,\k) circle (1.1pt);
  \draw[redacc, dashed, thick] (3,3.5) -- (3,14);
  \node[font=\tiny\bfseries, redacc, above] at (3,14) {observed 3};
  \node[font=\tiny] at (0,-7.5) {Simulated difference in mean yield (kg)};
  \node[font=\scriptsize\bfseries] at (0,33) {100 random reshuffles};
\end{tikzpicture}
\end{minipage}

\vspace{5pt}
\boxguard[10]
\textbf{Move 16 --- Scope of inference: two separate questions (3.5).} Was the treatment
\emph{assigned} at random? That gives \textbf{cause}. Were the individuals \emph{selected} at
random? That lets you \textbf{generalize} to the population they came from.

\vspace{3pt}
{\centering\footnotesize
\renewcommand{\arraystretch}{1.25}
\begin{tabular}{|c|c|c|}
\hline
 & \cellcolor{sky}\textbf{Selected at random} & \cellcolor{sky}\textbf{Not selected at random} \\
\hline
\cellcolor{sky}\textbf{Assigned at random} & Cause: yes \quad Generalize: yes & Cause: yes \quad Generalize: no \\
\hline
\cellcolor{sky}\textbf{Not assigned} & Cause: no \quad Generalize: yes & Cause: no \quad Generalize: no \\
\hline
\end{tabular}\par}
\normalsize

\vspace{0.10in}
\boxguard[22]
\begin{scenariobox}[The sentences that earn full credit]{navy}
\small
\textbf{Compare two distributions:} ``Weekend trips are typically longer (median 31 minutes
compared with 22) and more variable (IQR 14 minutes compared with 8). Weekday trips have one
high outlier, at 70 minutes.'' Center, variability, unusual features --- with comparison words
and numbers.\par\vspace{3pt}
\textbf{Outlier rule:} ``$Q_3 + 1.5(\text{IQR}) = 30 + 18 = 48$. Since $70 > 48$, the 70-minute
trip is an outlier.''\par\vspace{3pt}
\textbf{Slope:} ``For each additional hour of weekly practice, the \emph{predicted} recital score
increases by about 3.15 points.''\par\vspace{3pt}
\textbf{Residual:} ``This student scored 2.6 points \emph{higher} than the line predicted for 8
hours of practice.''\par\vspace{3pt}
\textbf{R-Sq:} ``About 81\% of the variation in recital score is explained by the linear model
with weekly practice hours.''\par\vspace{3pt}
\textbf{Bias and direction:} ``Nonresponse bias. Riders who rarely use the program are less
likely to answer, so they are under-represented, and the estimated proportion who ride weekly is
likely too high.''\par\vspace{3pt}
\textbf{Random assignment:} ``Label the plants 1 to 60. Randomly generate 30 different integers
from 1 to 60; those plants get mulch and the other 30 do not. This balances other variables, so
a difference in yield can be attributed to the mulch.''\par\vspace{3pt}
\textbf{Cause and to whom:} ``The treatments were randomly assigned and the difference is
significant, so the mulch caused the higher yield. The plants were not randomly selected, so the
result applies only to plants like these.''
\end{scenariobox}

\sghead{Ten traps that cost points}

\small
\begin{enumerate}[leftmargin=1.6em, itemsep=2pt]
  \item \textbf{Digits are not quantities.} A ZIP code is categorical.
  \item \textbf{Counts mislead when groups differ in size.} Compare proportions.
  \item \textbf{A longer whisker is not more data.} Every section of a boxplot holds a quarter.
  \item \textbf{The mean follows the tail; the median does not.}
  \item \textbf{The denominator is the group after ``of the.''} Read the question twice.
  \item \textbf{Residual $=$ actual $-$ predicted.} Negative means the line predicted too high.
  \item \textbf{$r$ takes the sign of the slope,} and a high $r$ does not prove a line fits.
  \item \textbf{A bigger sample does not fix bias.}
  \item \textbf{Blocking does not replace chance.} Randomize inside every block.
  \item \textbf{A random sample is not a random assignment.} Selection lets you generalize; only
        assignment lets you say cause.
\end{enumerate}
\normalsize

\sghead[6]{Before you walk in}

\small
Work the \textbf{practice final} in one sitting, timed, without notes, then check it against its
key. For anything you miss, go back to that unit's study guide and its sample test. A
free-response answer shorter than the key's is usually missing the context.
\normalsize

\end{document}
